\section{从“完蛋了”到客户倒戈}

本章进入创业低谷和拐点。视频里有一段“想说完蛋了”的章节：手上有一点激光技术，却没有足够激动人心的市场。这个时刻很关键，因为它迫使团队从科技创业者转向真正的市场创业者：不再只问技术能做什么，而是问客户为什么买、什么时候买、用什么价格买。

\begin{importantbox}{硬科技创业的第一次转向}
从“我有技术”到“客户为什么必须买”，是硬科技创业的关键转向。前者是研发逻辑，后者是产业逻辑。
\end{importantbox}

\subsection{定价心思：成本、价值和战略位置}

上一节说硬科技创业要从技术逻辑转向客户逻辑，本节把这个转向落到定价上。硬件定价不是成本加成。它要考虑 BOM 成本、客户安全价值、竞品替代价格、未来降本预期、进入主机厂的战略位置。早期价格定错，可能丢失客户；价格太低，又可能无法支撑交付和迭代。

\begin{figure}[H]
\centering
\includegraphics[width=0.96\textwidth]{pricing-logic.png}
\caption{硬件定价心思：定价不是成本加成，而是客户价值、风险和市场位置。自制概念图，依据 01:20:47--01:38:15 对谈内容整理。}
\end{figure}

\begin{knowledgebox}{读图：价格是战略语言}
硬件价格同时告诉客户三件事：这东西值多少钱、供应商承担多少风险、未来是否能规模化。早期定价如果只按成本算，会低估进入主机厂体系的战略价值。
\end{knowledgebox}

\begin{knowledgebox}{硬件定价的五个变量}
\begin{tabularx}{\textwidth}{p{0.20\textwidth}p{0.34\textwidth}X}
\toprule
变量 & 含义 & 影响 \\
\midrule
成本 & 物料、制造、良率、售后 & 决定底线。 \\
价值 & 对安全、体验、性能的提升 & 决定客户愿付价格。 \\
竞品 & 旧供应商和替代方案 & 决定切入空间。 \\
规模 & 未来量产降本空间 & 决定长期报价策略。 \\
战略 & 是否进入关键客户体系 & 决定短期利润是否可以让渡。 \\
\bottomrule
\end{tabularx}
\end{knowledgebox}

\begin{warningbox}{硬件低价策略的陷阱}
低价能加速客户倒戈，但也可能让公司长期承受负毛利、售后和质量压力。硬件价格必须同时服务进入市场和维持组织生命，不能只追求“拿下客户”。
\end{warningbox}

\subsection{客户开始倒戈}

定价说明供应商如何表达价值，本节看客户如何真正投票。客户从旧供应商转向新供应商，是硬件产业的突破时刻。它不是靠一次演示，而是靠成本、性能、可靠性、交付能力和行业信号共同作用。倒戈意味着客户相信新方案不只是便宜，而是足够可信。

\begin{figure}[H]
\centering
\includegraphics[width=0.96\textwidth]{customer-defection.png}
\caption{客户开始倒戈：硬件供应链的突破时刻，是客户愿意从旧方案转向新方案。自制概念图，依据 01:38:15--01:58:07 对谈内容整理。}
\end{figure}

\begin{knowledgebox}{客户为什么会倒戈}
\begin{tabularx}{\textwidth}{p{0.20\textwidth}p{0.34\textwidth}X}
\toprule
条件 & 客户看到什么 & 对供应商意味着什么 \\
\midrule
性能可信 & 能满足关键场景 & 技术不再只是 demo。 \\
成本优势 & 总成本明显下降 & 有机会进入量产车型。 \\
交付稳定 & 能按计划供货 & 组织能力被验证。 \\
替代风险可控 & 旧方案迁移不至于失控 & 客户愿意承担切换风险。 \\
\bottomrule
\end{tabularx}
\end{knowledgebox}

\subsection{本章小结}

从“完蛋了”到客户倒戈，关键是从技术自信转向客户验证。硬科技公司真正的拐点，是客户愿意把风险交给你。

